(1) 取质量 $m_2$ 向下的加速度 $a>0$，两端分别满足 $m_2g-T_2=m_2a$、$T_1-m_1g=m_1a$。无滑动给角加速度 $\alpha=a/R$；滑轮的转动方程是 $(T_2-T_1)R=I\alpha$，说明 $I>0$ 时两侧张力必不相等。

(2) 三式消元得 $\boxed{a=(m_2-m_1)g/(m_1+m_2+I/R^2)}$。张力为 $T_1=m_1(g+a)$、$T_2=m_2(g-a)$，两式均正且 $T_2>T_1$。由 $v^2=2ah$ 可得下降 $h$ 后 $v=\sqrt{2ah}$，滑轮角速度为 $v/R$。

(3) 重力势下降为 $(m_2-m_1)gh$，最终机械动能为 $\frac12(m_1+m_2)v^2+\frac12I(v/R)^2$。能量相等给 $v^2=2(m_2-m_1)gh/(m_1+m_2+I/R^2)$，与第二问独立一致。$I\to0$ 时恢复理想Atwood加速度；$I/R^2\gg m_1+m_2$ 时加速度极小。

(4) 若向下正转时轴承反扭矩为 $\tau_f$，转动方程变为 $(T_2-T_1)R-\tau_f=Ia/R$，故 $\boxed{a=[(m_2-m_1)g-\tau_f/R]/(m_1+m_2+I/R^2)}$。若静摩擦上限也为 $\tau_f$，从静止启动要求 $(m_2-m_1)gR>\tau_f$。下降 $h$ 的轴承热耗散为 $Q=\tau_f h/R$。
